\chapter{Mở Tiết Bằng Cử Chỉ}
\chapterintro{Không phải khởi động nào cũng cần lời}{Nhiều lớp vào nhịp nhanh hơn khi được làm một động tác rõ ràng trước khi phải nói.}

\section{Đứng nếu đã sẵn sàng}
\begin{khoidongbox}{Hoạt động khởi động}
Mời những em đã sẵn sàng đứng lên trong ba giây rồi ngồi xuống, không cần phát biểu gì thêm.
\end{khoidongbox}
\begin{visaohieuqua}
Động tác rất ngắn nhưng có tác dụng chuyển trạng thái cơ thể. Lớp thường tỉnh hơn sau một chuyển động đơn giản.
\end{visaohieuqua}
\begin{cachlam5phut}
Dùng đầu tiết sáng hoặc đầu giờ chiều. Sau đó gọi một em vừa đứng lên nói câu mở bài thật ngắn.
\end{cachlam5phut}
\ngancach

\section{Chỉ tay vào góc chọn đáp án}
\begin{khoidongbox}{Hoạt động khởi động}
Quy ước góc trái là \textit{đồng ý}, góc phải là \textit{còn phân vân}, rồi cho lớp chỉ tay chọn trong 5 giây.
\end{khoidongbox}
\begin{visaohieuqua}
Học sinh nhập cuộc cùng lúc mà không cần nói dài. Cử chỉ giúp lớp bật tham gia nhanh hơn gọi cá nhân từng em.
\end{visaohieuqua}
\begin{cachlam5phut}
Chỉ nên dùng với một câu hỏi rất nhẹ. Sau đó mời hai bạn ở hai phía giải thích thật ngắn.
\end{cachlam5phut}
\ngancach

\section{Vẽ biểu tượng bằng tay}
\begin{khoidongbox}{Hoạt động khởi động}
Yêu cầu lớp dùng ngón tay vẽ trên không một biểu tượng cho tiết học hôm nay: tia chớp, đám mây, hay mặt trời.
\end{khoidongbox}
\begin{visaohieuqua}
Hoạt động có hình ảnh khiến lớp thích hơn, nhất là những lớp còn ngại nói ngay đầu giờ. Nó vừa vui vừa ngắn.
\end{visaohieuqua}
\begin{cachlam5phut}
Sau khi cả lớp vẽ, hỏi 2 em vì sao chọn biểu tượng đó rồi vào bài bằng chính hình ảnh ấy.
\end{cachlam5phut}
\ngancach

\section{Giơ ngón tay chọn mức hiểu}
\begin{khoidongbox}{Hoạt động khởi động}
Nếu đây là tiết nối tiếp, mời lớp giơ 1, 2, hoặc 3 ngón tay để báo mức nhớ bài cũ.
\end{khoidongbox}
\begin{visaohieuqua}
Tín hiệu đơn giản cho thầy cô dữ liệu thật rất nhanh. Học sinh cũng thấy mình được hỏi trước khi bị kiểm tra.
\end{visaohieuqua}
\begin{cachlam5phut}
Nhìn mức số đông rồi quyết định ôn nhanh hay đi tiếp. Gọi đúng một em giải thích vì sao mình chọn mức đó.
\end{cachlam5phut}
\ngancach

\section{Cú xoay người đổi trạng thái}
\begin{khoidongbox}{Hoạt động khởi động}
Cho lớp xoay người sang bạn bên cạnh, nói một câu, rồi quay lại bảng trong chưa đến 20 giây.
\end{khoidongbox}
\begin{visaohieuqua}
Chỉ một cú xoay người cũng đủ tạo cảm giác tiết học đã bắt đầu. Cơ thể đổi hướng thì sự chú ý cũng đổi theo.
\end{visaohieuqua}
\begin{cachlam5phut}
Câu nói nên rất ngắn như: \textit{``Điều em nhớ nhất là gì?''} hoặc \textit{``Em đang ở mức mấy?''}
\end{cachlam5phut}